\section{Introduction}

To satisfy the rapidly growing computation demands of deep neural networks (DNNs), hardware vendors have introduced specialized floating-point matrix multiplication accelerators (MMAs) in their latest GPU architectures. MMAs such as NVIDIA Tensor Cores \cite{markidis_nvidia_2018,sun_dissecting_2023} and AMD Matrix Cores \cite{schieffer_rise_2024} empower modern GPUs with PetaFLOPS-scale computing capability, and play a pivotal role in the training and inference of frontier DNNs such as large language models (LLMs).

However, the rapid evolution of these high-performance MMAs introduces new challenges because no publicly available document specifies the exact arithmetic for floating-point matrix multiplication (\S\ref{sec:motivation}).  
Recent studies \cite{deepseek-ai_deepseek-v3_2024,pytorch_developers_pytorch_2025} have shown that insufficient arithmetic precision and dynamic range in some MMAs can lead to unexpected degradation in DNN training accuracy, particularly for LLMs. Diagnosing and preventing this issue is challenging because the internal arithmetic behaviors of MMAs are typically undocumented and thus poorly understood.

Moreover, different vendors adopt diverse approaches for the same matrix multiplication operations, and diversity exists even between different generations of MMAs from the same vendor \cite{li_discovery_2024}. These variations lead to subtle, sometimes significant numerical inconsistencies in matrix multiplication results, and consequently compromise the reproducibility of higher-level DNN training and inference.

To address these challenges, this paper presents MMA-Sim, the first bit-accurate reference model that reveals the detailed arithmetic behaviors of matrix multiplication accelerators (MMAs). Given input matrices with specified data types and shapes, MMA-Sim produces an output matrix that matches the MMA’s output bit by bit. By faithfully modeling each MMA’s arithmetic behavior, MMA-Sim enables thorough understanding, transparent analysis, and bit-accurate simulation of the subtle arithmetic intricacies of the MMA. With MMA-Sim, software developers can diagnose and prevent numerical issues in precision-sensitive DNN training and inference, and hardware designers can design new MMAs with numerical consistency to established and widely-adopted MMAs.

Bit-accurate modeling of MMAs is non-trivial. Most MMAs are black boxes for the public, meaning that their arithmetic implementations are not disclosed. To uncover the detailed arithmetic behaviors of each matrix multiplication instruction of each MMA, we employ specially crafted inputs along with randomized inputs to detect the characteristics of the MMA in both edge and typical cases (\S\ref{sec:method}). Based on the observed results, we construct an arithmetic algorithm that models the arithmetic behavior of the instruction. The algorithm specifies the step-by-step computation procedure, including the operation type, operation order, rounding method, precision, and other implementation details. The correctness of each algorithm is carefully validated against the MMA output at the bit level, and iteratively refined until bitwise equivalence is achieved.

Applying this testing-based approach to all matrix multiplication instructions from ten GPU architectures, we construct nine arithmetic algorithms (\S\ref{sec:model}). Using the nine algorithms, we implement MMA-Sim using approximately 2800 lines of Python code. MMA-Sim faithfully simulates the diverse arithmetic of all floating-point matrix multiplication instructions across all the eight NVIDIA GPU architectures with Tensor Cores (from Volta to RTX Blackwell) and the two latest AMD GPU architectures with Matrix Cores (CDNA2 and CDNA3). MMA-Sim is validated against real hardware through extensive large-scale random testing (\S\ref{sec:correctness}). We test each instruction with more than one million sets of randomly generated input matrices, and the results confirm bitwise equivalence between MMA-Sim and the corresponding hardware.

Through MMA-Sim, we investigate previously reported issues and uncover many previously unknown characteristics in MMA designs across vendors and architectures (\S\ref{sec:findings}). For example, we confirm the reduced accumulation precision of the 8-bit floating-point (FP8) instructions on the NVIDIA Hopper architecture, and the reduced dynamic range of the 16-bit floating-point (FP16 and BF16) instructions on the AMD CDNA2 architecture, both of which can adversely affect DNN training, as reported in previous work \cite{deepseek-ai_deepseek-v3_2024,pytorch_developers_pytorch_2025}. We also discover the reduced accumulation precision of the FP8 instructions on the NVIDIA Ada Lovelace architecture, and find that the precision is increased on the newer Blackwell and RTX Blackwell architectures. Moreover, we identify the asymmetric rounding used in the TF32 (TensorFloat-32), FP16, and BF16 instructions on the AMD CDNA3 architecture, which can introduce notable numerical errors and biases in certain conditions. Interestingly, we also find that the rounding method has been specially modified in the FP8 instructions on the CDNA3 architecture, which mitigates such issues.

These findings demonstrate the power of MMA-Sim as an analytical tool to understand the detailed arithmetic behaviors of modern accelerators. For software developers, MMA-Sim bridges the gap between hardware implementations and software expectations, and enables systematic investigation of numerical error sources, stability analysis of training algorithms, and hardware-specific precision tuning. For hardware designers, MMA-Sim as a bit-accurate reference model captures the full spectrum of arithmetic subtleties, providing the detailed functional specifications for future MMA design and verification.


In summary, this paper makes the following contributions.

\begin{itemize}
\item Systematically dissecting MMA arithmetic behaviors. We develop a testing-based methodology that detects each MMA’s arithmetic behavior through specially crafted inputs and random inputs, and derive nine distinct arithmetic algorithms that capture their arithmetic for floating-point matrix multiplication.
\item Bit-accurate reference model. We present MMA-Sim, the first bit-accurate reference model that simulates the exact arithmetic behaviors of modern MMAs across vendors and architectures with the nine arithmetic algorithms.
\item Implementation and validation. We implement MMA-Sim for ten GPU architectures (eight from NVIDIA and two from AMD) and validate its correctness via large-scale random testing, confirming bitwise equivalence with the real hardware.
\item New insights into MMA arithmetic design and usage. Through MMA-Sim, we identify many undocumented design choices that explain the reported instability in DNN training or introduce new numerical issues, and we suggest mitigation approaches to both hardware and software developers.
\item Practical analytical utility. MMA-Sim will be open-source, providing a unified framework for MMA arithmetic behavior modeling, numerical analysis, arithmetic design specification, and hardware-aware optimization in both software and hardware design.
\end{itemize}
